% Generated draft: requires mathematical and source review.
\documentclass[11pt]{article}
\usepackage[margin=25mm]{geometry}
\usepackage{fontspec}
\setmainfont{Latin Modern Roman}
\usepackage{amsmath,amssymb,mathtools,unicode-math}
\setmathfont{Latin Modern Math}
\usepackage{enumitem,needspace,xurl,hyperref}
\setlength{\parindent}{0pt}
\setlength{\parskip}{5pt}
\setlength{\emergencystretch}{3em}
\newcommand{\N}{\mathbb{N}}
\newcommand{\Z}{\mathbb{Z}}
\newcommand{\Q}{\mathbb{Q}}
\newcommand{\R}{\mathbb{R}}
\newcommand{\C}{\mathbb{C}}
\newcommand{\F}{\mathbb{F}}
\newcommand{\A}{\mathbb{A}}
\newcommand{\PP}{\mathbb{P}}
\newcommand{\E}{\mathbb{E}}
\newcommand{\eps}{\varepsilon}
\newcommand{\dd}{\,\mathrm{d}}
\begin{document}
\section*{UnsolvedMath ID: 30002447}
\textbf{Source problem code:} catalogue entry
\subsection{Definitions and mathematical statement}
The supplied question concerns finite-difference and finite-element approximations of PDE eigenvalue problems. Neither a particular PDE eigenvalue problem nor a discretization is specified.

A conditional counting framework is as follows. After fixing a PDE eigenvalue problem, a discretization, and its parameters, write its finite list of computed eigenvalues as \((\widehat\lambda_j)_{j\in J}\), where \(J\) is a finite index set. Once a reliability criterion has been specified, let \(R(j)\) mean that the computed eigenvalue with index \(j\) satisfies that criterion, and define
\[
N_R:=\#\{j\in J:R(j)\}.
\]
This definition counts entries of the computed list; whether the intended question counts distinct eigenvalues or eigenvalues with multiplicity is unspecified.

\textbf{Source question.} How many numerical eigenvalues from finite-difference or finite-element approximations of PDE eigenvalue problems are reliable?

In the conditional framework above, the requested quantity would be \(N_R\), or quantitative bounds for it in terms of the chosen problem and discretization parameters. However, the source does not define \(R\), specify an accuracy tolerance or convergence requirement, or give the operator, domain, boundary conditions, spectral assumptions, or approximation hypotheses. It also does not specify whether the desired result is an asymptotic count, a finite-resolution estimate, or a certified count. Consequently, the supplied question is a research direction rather than a fully quantified mathematical assertion; the framework above does not supply the missing hypotheses or select a reliability criterion.

The supplied dated curation labels the problem partially solved and reports that quantitative trusted-eigenvalue counts are known for standard elliptic finite-element and finite-difference models, while no method-independent answer is available and spectral pollution remains possible. These are supplied literature qualifications, not independently verified conclusions of this draft.
\subsection{Short English statement}
The question asks how many computed PDE eigenvalues can be trusted when using finite differences or finite elements. A precise count requires a specified PDE problem, numerical method, and meaning of reliability. Those choices are absent from the supplied record. Its dated literature curation reports partial results for standard elliptic models and warns that spectral pollution remains possible.
\subsection{Sources}
\begin{itemize}
\item Supplied local UnsolvedMath registry record, permanent ID \(30002447\), title \emph{Reliability of Numerically Computed PDE Eigenvalues}, source problem code \texttt{catalogue entry}, in the collection \emph{Oberwolfach Reports Open Problems}. No individual proposer is identified. The record gives proposed year 2013.
\item Supplied report citation: \emph{Numerical Solution of PDE Eigenvalue Problems} (2014), Oberwolfach Reports, pp.~3221--3304, DOI \texttt{10.4171/owr/2013/56}. Supplied literature-source label: original OWR report. \url{https://doi.org/10.4171/owr/2013/56}. The background describes Oberwolfach Reports as documenting mathematical workshops and their open-problem discussions.
\item Supplied literature-source label: Zhang trusted eigenvalues. \url{https://doi.org/10.1007/s10915-014-9971-5}. Full bibliographic details and a specific supporting theorem are not supplied.
\item Supplied literature-source label: spectral pollution and bounds. \url{https://doi.org/10.1016/j.apnum.2015.08.007}. Full bibliographic details and a specific supporting theorem are not supplied.
\item Literature assessment date: \texttt{2026-08-21}. This is supplied dated curation, described in the record as literature triage derived from the supplied OWR curation, and has not been independently verified. The supplied status is \texttt{partially\_solved}. The assessment reports quantitative trusted-eigenvalue counts for standard elliptic finite-element and finite-difference models, no method-independent answer, and the possibility of spectral pollution. Time-sensitive status and cited claims require independent verification.
\item The registry marks the extracted statement as exact and reports manual review for self-containment, mathematical coherence, and absence of LaTeX or report-layout artifacts. These are supplied verification claims; they do not resolve the missing mathematical specifications identified here.
\end{itemize}
\end{document}
